\chapter[Noise and Error]{Noise, Error, and the Limits of Real Experiments}
\label{chap:noise-error-and-the-limits-of-real-experiments}

\begin{learningobjectives}
\begin{itemize}
  \item Interpret the symmetric outcome-flip model.
  \item Distinguish readout error, bit flips, and phase errors conceptually.
  \item Compare error mitigation with error correction.
  \item Explain the purpose and limitation of a repetition code.
\end{itemize}
\end{learningobjectives}

\begin{openingquestion}
How can we distinguish sampling fluctuation from a model of device error?
\end{openingquestion}

\paragraph{App and experiment mapping.}
Primary app material: \applevel{8}, \applevel{10}, \applevel{13}. Reproducible experiment IDs:
\experiment{EXP-N01, EXP-STAT02, EXP-EC01, EXP-EC02}.

\section{Try it first}
\begin{tryit}
\placeholder{Insert a short app procedure using the experiment registry. Include
the initial settings, what the learner changes, and what should be recorded before
the explanation is revealed.}
\end{tryit}

\section{What happened?}
\placeholder{Describe the visible result without yet introducing unnecessary formalism.}

\section{The main idea}
\begin{coreidea}
\placeholder{State one central claim in plain language.}
\end{coreidea}

\section{The mathematics}
\placeholder{Introduce only the mathematics needed for this chapter, defining every symbol.}

\section{Connecting the mathematics and the code}
\placeholder{Map mathematical objects to tested simulator functions and generated figures.
Do not paste the full Streamlit page implementation into the manuscript.}

\section{Common misconceptions}
\begin{misconception}
\placeholder{State a common incorrect interpretation, explain why it fails, and provide
a counterexample from the app.}
\end{misconception}

\section{Check your understanding}
\begin{checkpoint}
\begin{enumerate}
  \item \placeholder{Prediction question}
  \item \placeholder{Short calculation}
  \item \placeholder{Interpretation of a graph or circuit}
\end{enumerate}
\end{checkpoint}

\section{Going deeper}
\begin{goingdeeper}
\placeholder{Optional mathematical treatment. The core chapter must remain readable if
this box is skipped.}
\end{goingdeeper}

\section{Further reading}
\placeholder{Add primary or authoritative references after the literature matrix is built.}
